\begin{definitionbox}{Definition --- Basis for a Topology}
\begin{definition}[Basis for a Topology]\label{def:topological-basis}
Let \((X,\tau)\) be a topological space.  A collection
\(\mathcal{B}\subseteq\tau\) is a \emph{basis for the topology} \(\tau\) if
every open set \(O\in\tau\) is the union of some subcollection of
\(\mathcal{B}\).  The elements of \(\mathcal{B}\) are called
\emph{basic elements} or \emph{basic open sets}.
\end{definition}
\end{definitionbox}

\begin{remark*}[Standard quantified statement]
\[
\mathcal{B}\text{ is a basis for }\tau
\Longleftrightarrow
\mathcal{B}\subseteq\tau
\text{ and }
\forall O\in\tau\,\exists\mathcal{A}\subseteq\mathcal{B}{:}\;
O=\bigcup_{B\in\mathcal{A}}B.
\]
\end{remark*}

\begin{remark*}[Interpretation]
A basis is a smaller supply of open sets from which all open sets can be
assembled by unions.  It records local building blocks for the topology.
\end{remark*}

\begin{dependencies}
\begin{itemize}
  \item \hyperref[def:topological-space]{Definition: Topological Space}
\end{itemize}
\end{dependencies}

\begin{propositionbox}{Proposition}
\begin{proposition}[Pointwise Characterization of a Basis]
\label{prop:topological-basis-pointwise-characterization}
Let \((X,\tau)\) be a topological space, and let
\(\mathcal{B}\subseteq\tau\).  Then \(\mathcal{B}\) is a basis for \(\tau\)
if and only if for every open set \(O\in\tau\) and every \(x\in O\), there is
\(B\in\mathcal{B}\) such that
\[
  x\in B\subseteq O.
\]
Consequently, a subset \(A\subseteq X\) is open if and only if for each
\(x\in A\), there is \(B\in\mathcal{B}\) such that
\[
  x\in B\subseteq A.
\]

\smallskip
\noindent
\hyperref[prf:topological-basis-pointwise-characterization]{\textit{Go to proof.}}
\end{proposition}
\end{propositionbox}

\begin{remark*}[Standard quantified statement]
\[
\mathcal{B}\text{ is a basis for }\tau
\Longleftrightarrow
\forall O\in\tau\,\forall x\in O\,\exists B\in\mathcal{B}{:}\;
x\in B\subseteq O.
\]
\[
A\in\tau
\Longleftrightarrow
A\subseteq X
\text{ and }
\forall x\in A\,\exists B\in\mathcal{B}{:}\;
x\in B\subseteq A.
\]
\end{remark*}

\begin{remark*}[Interpretation]
The pointwise test says that a basis can be checked locally.  Around every
point of every open set, there must be a basic open set that still lies inside
the original open set.
\end{remark*}

\begin{dependencies}
\begin{itemize}
  \item \hyperref[def:topological-space]{Definition: Topological Space}
  \item \hyperref[def:topological-basis]{Definition: Basis for a Topology}
\end{itemize}
\end{dependencies}
